\chapter{Extremal Problems and Ramsey Theory}
\label{ch:extremal}

\begin{center}\itshape
Among graphs on $n$ vertices with no clique of size $r+1$, the Tur\'an graph has the most edges.\\
Every large enough colouring of a complete graph contains a monochromatic clique.
\end{center}

The arguments below are counting arguments. They do not use a later volume.

\section{Cliques and colourings}

The Tur\'an graph $T(n,r)$ is the complete $r$-partite graph on $n$ vertices with parts as equal as possible.

\begin{theorem}[Tur\'an]\label{thm:turan}
Among $K_{r+1}$-free graphs on $n$ vertices, $T(n,r)$ has the greatest number of edges.
\end{theorem}
\begin{proof}
A maximum graph is complete multipartite: if two nonadjacent vertices had different neighbourhoods, copying the larger neighbourhood onto the other vertex would add edges without creating a $K_{r+1}$. Nonadjacency is then an equivalence relation, and there are at most $r$ classes. Balancing the class sizes increases the number of edges whenever two classes differ by two or more.
\end{proof}

\begin{theorem}[Ramsey]\label{thm:ramsey}
The least $R(s,t)$ such that every red-blue colouring of $K_{R(s,t)}$ contains a red $K_s$ or a blue $K_t$ is finite, and $R(s,t)\le R(s-1,t)+R(s,t-1)$. Also $R(3,3)=6$.
\end{theorem}
\begin{proof}
In a colouring of a complete graph on $R(s-1,t)+R(s,t-1)$ vertices, some vertex has red degree at least $R(s-1,t)$ or blue degree at least $R(s,t-1)$. A red clique in the red neighbourhood, or a blue clique already there, finishes the induction. The values $R(2,t)=t$ and $R(s,2)=s$ start it. For the upper bound $R(3,3)\le6$, a vertex of $K_6$ has three edges of one colour; those three neighbours close a monochromatic triangle. For the lower bound, a red $5$-cycle and its blue complement have no monochromatic triangle.
\end{proof}

\begin{theorem}[Schur]\label{thm:schur}
For every $r\ge1$ there exists $S$ such that every $r$-colouring of $\{1,\ldots,S\}$ contains monochromatic $x,y,z$ with $x+y=z$.
\end{theorem}
\begin{proof}
First, the $r$-colour Ramsey number for a monochromatic triangle is finite. For $r=1$ it is $3$. If it is finite for $r-1$, call it $N$. In an $r$-edge-colouring of a complete graph on $r(N-1)+2$ vertices, some vertex has at least $N$ edges of one colour, say red. If any edge inside that neighbourhood is red, those three vertices form a red triangle. If none is red, the neighbourhood is coloured with the remaining $r-1$ colours and contains a monochromatic triangle. Colour the edge $\{i,j\}$ of this complete graph by the colour of the positive integer $|i-j|$. A monochromatic triangle $a<b<c$ yields $x=b-a$, $y=c-b$ and $z=c-a$, all the same colour, with $x+y=z$.
\end{proof}

\section{Sunflowers and intersecting sets}

A sunflower with $k$ petals is a collection of $k$ sets with all pairwise intersections equal.

\begin{theorem}[Erd\H os--Rado]\label{thm:sunflower}
More than $s!(k-1)^s$ sets of size $s$ contain a sunflower with $k$ petals.
\end{theorem}
\begin{proof}
Induct on $s$. Distinct singletons give the case $s=1$. If some element lies in more than $(s-1)!(k-1)^{s-1}$ of the sets, delete it and apply the inductive hypothesis. Otherwise a maximal disjoint subfamily has at least $k$ members: fewer than $k$ disjoint sets would cover the family by at most $(k-1)s(s-1)!(k-1)^{s-1}$ sets. A disjoint family is a sunflower with empty core.
\end{proof}

\begin{theorem}[Erd\H os--Ko--Rado]\label{thm:ekr}
If $n\ge 2k\ge2$, an intersecting family of $k$-subsets of an $n$-set has size at most $\binom{n-1}{k-1}$.
\end{theorem}
\begin{proof}
There are $(n-1)!$ ways to place the $n$ points on a directed cycle, up to rotation. On each cycle at most $k$ of the arcs of length $k$ are pairwise intersecting, because two arcs that begin at least $k$ steps apart are disjoint when $n\ge2k$. Each $k$-set lies on $k!(n-k)!$ directed cycles. Counting pairs of a cycle and a set of the family on that cycle gives $|\mathcal F|\,k!(n-k)!\le(n-1)!\,k$.
\end{proof}

The family of all $k$-sets through a fixed point meets the bound.

\section{Exercises}

\begin{exercise}[$\star$]\label{ex:r33}
Give a $2$-edge-colouring of $K_5$ with no monochromatic triangle, and the counting that forces one in $K_6$.
\end{exercise}

\begin{exercise}[$\star\star$]\label{ex:turanedges}
Show that $T(n,2)$ has $\lfloor n^2/4\rfloor$ edges.
\end{exercise}

\begin{exercise}[$\star\star$]\label{ex:ekrsharp}
Show that the family of all $k$-subsets containing a fixed element is intersecting and has size $\binom{n-1}{k-1}$.
\end{exercise}

\section{Hints}

\begin{hint}{ex:r33}
Use a $5$-cycle.
\end{hint}
\begin{hint}{ex:turanedges}
Multiply the two part sizes.
\end{hint}
\begin{hint}{ex:ekrsharp}
Choose the remaining $k-1$ elements from the other $n-1$.
\end{hint}

\section{Solutions}

\begin{solution}{ex:r33}
\writeup A red $5$-cycle and a blue complementary $5$-cycle have no monochromatic triangle. In $K_6$ some vertex has three edges of one colour, and those three neighbours complete a monochromatic triangle.
\end{solution}

\begin{solution}{ex:turanedges}
\writeup The parts have sizes $\lfloor n/2\rfloor$ and $\lceil n/2\rceil$, and every edge between them is present, so there are $\lfloor n^2/4\rfloor$ edges.
\end{solution}

\begin{solution}{ex:ekrsharp}
\writeup Any two sets that both contain the fixed element intersect there. The count is the number of ways to finish the set from the remaining $n-1$ elements.
\end{solution}

\section*{Sources and further reading}

\begin{enumerate}
\item P.~Tur\'an, Eine Extremalaufgabe aus der Graphentheorie, \emph{Mat.\ Fiz.\ Lapok} \textbf{48} (1941), 436--452.
\item F.~P.~Ramsey, On a problem of formal logic, \emph{Proc.\ London Math.\ Soc.} \textbf{30} (1930), 264--286.
\item P.~Erd\H os, C.~Ko and R.~Rado, Intersection theorems for systems of finite sets, \emph{Quart.\ J.\ Math.\ Oxford} \textbf{12} (1961), 313--320.
\item G.~O.~H.~Katona, A simple proof of the Erd\H os--Ko--Rado theorem, \emph{J.\ Combin.\ Theory Ser.\ B} \textbf{13} (1972), 183--184.
\item P.~Erd\H os and R.~Rado, Intersection theorems for systems of sets, \emph{J.\ London Math.\ Soc.} \textbf{35} (1960), 85--90. The sunflower lemma.
\end{enumerate}
